\documentclass[11pt,a4paper]{article}
\usepackage[T1]{fontenc}
\usepackage{lmodern,amsmath,amssymb,amsthm,mathtools,booktabs,microtype}
\usepackage[margin=27mm]{geometry}
\usepackage[colorlinks=true,linkcolor=blue,citecolor=blue,urlcolor=blue]{hyperref}
\setlength{\emergencystretch}{2em}
\newtheorem{theorem}{Theorem}[section]
\newtheorem{lemma}[theorem]{Lemma}
\newtheorem{corollary}[theorem]{Corollary}
\newtheorem{proposition}[theorem]{Proposition}
\newcommand{\Q}{\mathbb Q}
\newcommand{\Z}{\mathbb Z}
\title{A local obstruction to affine identifications of scalar and input points}
\author{}\date{16 November 2026}
\begin{document}\maketitle
\begin{abstract}
For the general elliptic detector with input polynomial
$216X^4+450459$ and its specified scalar coefficient, every fixed
affine identification $S=[m]A+[k]P$ of the scalar point with the
input point excludes an infinite arithmetic progression of integer
inputs. The same holds for every fixed finite union of such
identifications. The assertion holds throughout the source's sufficiently
large coefficient range. The proof combines its scalar local
inequality with an elementary prime-divisor argument. For
$S=A+[2]P$, the integer input $X=11$ gives an exact obstruction
at the prime $80287$. Consequently a simultaneous recovery of
$S$ and its auxiliary point by these affine formulas cannot remove
four rational coordinates while retaining completeness. The result
concerns the specified source and coefficient functions, not a
lower bound for Diophantine definitions.
\end{abstract}

\section{The source conditions}

Put
\[
 K=\Q(\sqrt2),\qquad
 E:\ y^2=x^3-(3-\sqrt2)^2x,
\]
and let
\[
 P=\left(\frac{57-28\sqrt2}{8},
         \frac{-660+417\sqrt2}{32}\right).
\]
We use the general guarded source of~\cite[Section~3]{Unit}, with
the original coefficient functions below. Its point roles include
$A,S,H,U_A,U_S$, with
\[
                         A=[n]P,\qquad S=[a]P,
                         \qquad n,a\in\Z.
\]
The specified source premise is $2E(K)=\Z P$, with $P$ of infinite
order. The source guards exclude $S-P=O$; in particular $a\ne1$.
All point and pole guards are retained.

For each fixed odd integer $L\ge3$, set
\begin{equation}\label{eq:coefficient}
\begin{gathered}
 \beta(X)=216X^4+450459,\qquad \ell=76249,\\
 r(X)=\beta(X)(\beta(X)-1)(\beta(X)-2401)
                         (32\beta(X)+343)(2\beta(X)-1),\\
 C_L=17(\ell^{(L+1)/2}+1),\qquad
 \rho_L(X)=\frac{102^2\ell^L r(X)}{(1+C_Lr(X))^2}.
\end{gathered}
\end{equation}
The exceptional set is
\begin{equation}\label{eq:exceptional}
 \mathcal S=\{p\text{ prime}:p\le76176\}
                    \cup\{1875611,3828397,3982747\}.
\end{equation}
In particular $\ell\notin\mathcal S$.

The obstruction uses exactly the following two source conclusions.
\begin{enumerate}
\item For every rational source solution in the sufficiently large
odd-$L$ range of~\cite[Section~3]{Unit},
\begin{equation}\label{eq:inputindex}
                              n=\beta(X).
\end{equation}
This is the vanishing of the input residual in that source's
soundness proof, not merely its conclusion that $X$ is integral.
\item For every odd $L\ge3$, every rational source solution and every
$p\notin\mathcal S$, if $\mu=v_p(\rho_L(X))$ is odd, then
$\mu>0$ and
\begin{equation}\label{eq:scalar}
                              v_p(a-1)>\frac94\mu.
\end{equation}
In particular $\mu=1$ forces $v_p(a-1)\ge3$.
\end{enumerate}
The scalar-odd lemma of~\cite[Section~3]{General} proves
\eqref{eq:scalar}. Section~3 of~\cite{Unit} verifies that its local
argument survives the unsquared comparison scale, and proves
\eqref{eq:inputindex}. Allowing the final pole ordinate to vanish
as in~\cite{Terminal} preserves these conclusions.

When applying these statements to the defining-polynomial family,
we retain all arithmetic premises of~\cite[Section~1.1]{General}
and~\cite[Section~1]{Unit}: the specified elliptic subgroup and
height, normalized modular realization and its contact-height
consequence, prepared-input pole statements, and simultaneous
prime-pattern input. No new arithmetic hypothesis is introduced.
The proof below is an elementary implication of
\eqref{eq:inputindex}--\eqref{eq:scalar}; it also applies to any
other guarded source retaining these two properties and
\eqref{eq:coefficient}--\eqref{eq:exceptional}.

\section{Prime divisors of polynomial values}

\begin{lemma}\label{lem:schur}
If $f\in\Z[T]$ is nonconstant, infinitely many primes divide its
nonzero values at integers.
\end{lemma}
\begin{proof}
Suppose all such primes belong to a finite set with product $M$,
where the empty product is one. Choose $n_0\in\Z$ with
$A=f(n_0)\ne0$. The polynomial
\[
                    g(T)=\frac{f(n_0+AMT)}{A}
\]
has integer coefficients, is nonconstant, and has the form
$1+MT h(T)$ with $h\in\Z[T]$. A sufficiently large integer
argument gives $|g(T)|>1$. Every prime divisor of this value
then lies outside the assumed set, since $g(T)\equiv1$ at
all its primes. Such a prime also divides the nonzero integer
$f(n_0+AMT)=A g(T)$, a contradiction.
\end{proof}

\begin{lemma}\label{lem:simplelift}
Let $f\in\Z[T]$ and let $p$ be a prime such that
$f(x_0)\equiv0\pmod p$ and $f'(x_0)\not\equiv0\pmod p$.
There is a residue class $c\pmod{p^2}$ such that
\[
                         v_p(f(x))=1
                         \qquad(x\equiv c\pmod{p^2}).
\]
\end{lemma}
\begin{proof}
For $j\in\Z$ the Taylor expansion gives
\[
 \frac{f(x_0+pj)}p
          \equiv\frac{f(x_0)}p+jf'(x_0)\pmod p.
\]
Exactly one class of $j$ makes this zero. Choose another class;
the resulting value of $x_0+pj$ modulo $p^2$ has the required
property. Polynomial values are unchanged modulo $p^2$ on that
class.
\end{proof}

\section{Exclusion of every fixed affine identification}

\begin{theorem}\label{thm:affine}
Fix integers $m,k$. If $(m,k)\ne(0,1)$, there are infinitely many primes
$p\notin\mathcal S\cup\{\ell\}$, each with a residue class
$c\pmod{p^2}$ such that, for every integer $X$ in that class and
every odd $L\ge3$, the relation
\begin{equation}\label{eq:affinelabel}
                              a=m\beta(X)+k
\end{equation}
is incompatible with the scalar local condition
\eqref{eq:scalar}. The prime and residue class are independent of
$L$. If $(m,k)=(0,1)$, the relation gives $a=1$, excluded by the
source guard.

Consequently, throughout the sufficiently large odd-$L$ range
where \eqref{eq:inputindex} holds, adjoining the point identity
\begin{equation}\label{eq:pointidentity}
                              S=[m]A+[k]P
\end{equation}
excludes infinitely many integer inputs from the guarded source.
\end{theorem}
\begin{proof}
Suppose $(m,k)\ne(0,1)$. Choose $t\in\{0,1\}$ by
\[
 t=0\quad\text{if }k\ne1,\qquad
 t=1\quad\text{if }k=1.
\]
Then
\begin{equation}\label{eq:nonzeroconstant}
                               mt+k-1\ne0.
\end{equation}
The nonconstant polynomial
\[
                    f_t(T)=216T^4+450459-t
\]
is squarefree over $\Q$: its derivative is $864T^3$, and
$f_t(0)=450459-t\ne0$. By Lemma~\ref{lem:schur}, infinitely many
primes divide nonzero values of $f_t$. Choose such a prime $p$
outside $\mathcal S\cup\{\ell\}$ and outside the finite prime
support of
\[
                        864(450459-t)(mt+k-1).
\]
A root $x_0$ of $f_t$ modulo this prime is nonzero modulo $p$,
and $f_t'(x_0)=864x_0^3$ is a unit. Lemma~\ref{lem:simplelift}
provides a residue class $c\pmod{p^2}$ on which
\begin{equation}\label{eq:valuationone}
                   v_p(\beta(X)-t)=1,\qquad
                   \beta(X)\equiv t\pmod p.
\end{equation}

If $t=0$, the four other factors of $r(X)$ reduce to
\[
                              -1,\quad-2401,\quad343,\quad-1.
\]
If $t=1$, they reduce to
\[
                              1,\quad-2400,\quad375,\quad1.
\]
All are units because $p>76176$. Thus
$v_p(r(X))=1$. Also $102\ell$ is a unit at $p$, whereas
$1+C_Lr(X)\equiv1\pmod p$ for every fixed odd $L\ge3$.
Equation~\eqref{eq:coefficient} therefore gives
\[
                               v_p(\rho_L(X))=1.
\]
The scalar condition requires $v_p(a-1)\ge3$. But
\eqref{eq:affinelabel}, \eqref{eq:valuationone}, and the choice
of $p$ give
\[
                    a-1\equiv mt+k-1\not\equiv0\pmod p.
\]
This is a contradiction for every integer $X$ in the residue class.
The prime-divisor lemma permits infinitely many choices of $p$ outside
the specified finite exclusions.

Finally, $P$ has infinite order, so
\eqref{eq:pointidentity} is equivalent to $a=mn+k$.
Using \eqref{eq:inputindex} gives \eqref{eq:affinelabel}.
This proves the point-identity conclusion in the stated coefficient
range. The pair $(m,k)=(0,1)$ gives $S=P$ and is excluded directly.
\end{proof}

\begin{corollary}\label{cor:finiteunion}
Let $\mathcal C\subset\Z^2$ be finite and nonempty. There is an
infinite arithmetic progression of integer inputs for which no source
solution satisfies $S=[m]A+[k]P$ for any $(m,k)\in\mathcal C$,
throughout the sufficiently large coefficient range. Thus a finite
union of these source restrictions is not complete.
\end{corollary}
\begin{proof}
The pair $(0,1)$ is excluded by the source guard for every input and
can be discarded. If no pair remains, the conclusion is immediate.
Otherwise choose a distinct prime and its obstructing residue class
from Theorem~\ref{thm:affine} for each remaining pair. This is possible
because each pair permits infinitely many primes. The Chinese
remainder theorem gives a simultaneous residue class modulo the
product of their squares. Every integer in this class has the local
obstruction for every allowed pair.
\end{proof}

\begin{corollary}\label{cor:two}
The simultaneous substitutions
\[
                  S=[m]A+[k]P,\qquad
                  U_S=[m]U_A+[k]H
\]
with fixed integers $m,k$ cannot preserve completeness of the
specified general source. In particular these substitutions cannot
remove the two $K$-point blocks, and hence four retained rational
coordinates, while keeping that source's integer input projection.
\end{corollary}
\begin{proof}
The first identity alone has the obstruction in
Theorem~\ref{thm:affine}. Adding the second identity cannot restore
an excluded input.
\end{proof}

The theorem concerns the coefficient functions
\eqref{eq:coefficient} and the two source properties
\eqref{eq:inputindex}--\eqref{eq:scalar}. It makes no assertion
about altered scalar coefficients, point identities involving other
root roles, or different Diophantine constructions.

\section{An exact counterexample for a two-step translation}

\begin{proposition}\label{prop:explicit}
For $S=A+[2]P$, one can take $X=11$ and $p=80287$.
The obstruction persists for all integers
\[
                              X\equiv11\pmod{6446002369}.
\]
\end{proposition}
\begin{proof}
Direct arithmetic gives
\[
 \beta(11)=3612915=3^2\cdot5\cdot80287,\qquad
                         80287^2=6446002369.
\]
The number $80287$ is prime, lies outside~\eqref{eq:exceptional},
and differs from $\ell$. The five factors of $r(11)$ have residues
\[
                    0,\quad80286,\quad77886,\quad343,\quad80286
                    \pmod{80287}.
\]
The first factor has valuation one, so
$v_{80287}(\rho_L(11))=1$ for every odd $L\ge3$.
For $a=\beta(11)+2$, however,
\[
                            a-1\equiv1\pmod{80287}.
\]
This contradicts~\eqref{eq:scalar}. The same factor residues and
first-order valuation hold whenever $X\equiv11\pmod{80287^2}$.
In the sufficiently large source range, the point identity forces
$a=\beta(X)+2$ by~\eqref{eq:inputindex}.
\end{proof}

\section*{Exact verification}
The accompanying standard-library script
\begin{center}\small
\texttt{negative\_results/computation/affine\_coupling\_v50/check.py}
\end{center}
checks primality by trial
division through $\lfloor\sqrt{80287}\rfloor$, the exact
factorization of $\beta(11)$, every contact residue, and the first-order
valuation. It writes \nolinkurl{affine_coupling_certificate.json}.
The proof of Theorem~\ref{thm:affine} gives the infinite-family
conclusion for every fixed pair $(m,k)$; it does not depend on a
bounded computational search.

\begin{thebibliography}{9}
\bibitem{General}
\emph{A defining polynomial in twenty-five variables: combined scalar
and comparison conditions}, repository source,
\nolinkurl{papers/few_variables25/few_variables25_merged_20261110_v19.tex}.
\bibitem{Unit}
\emph{A single norm determinant for the 25- and 13-variable constructions},
repository source,
\nolinkurl{papers/finite_algebra/unit_norm_20261114_v28.tex}.
\bibitem{Terminal}
\emph{Smaller norm polynomials with a possibly singular final quadratic
extension}, repository source,
\nolinkurl{papers/finite_algebra/terminal_quadratic_root_20261115_v40.tex}.
\end{thebibliography}
\end{document}
